\documentclass[12pt]{article}
\usepackage{amsmath, amssymb}
\usepackage[margin=1in]{geometry}

\begin{document}

\section*{Lecture 26: Stochastic Processes (SSS, WSS, Ergodic)}

\section*{1. Basics of Stochastic Processes}

\subsection*{Definition}

A \textbf{stochastic process} $X(\xi, t)$ is defined as:

\begin{itemize}
\item A \textbf{collection of waveforms/functions over time}
\item Each waveform depends on the outcome $(\xi)$ of a \textbf{random experiment}
\end{itemize}

\subsection*{Reasoning}

\begin{itemize}
\item A random experiment produces outcomes $\rightarrow$ each outcome generates one waveform.
\item Therefore, instead of one function, we get a \textbf{set of possible functions}.
\end{itemize}

\subsection*{Realizations and Random Variables}

\textbf{Key Concepts}

\begin{enumerate}
\item \textbf{Realization}

A \textbf{single observed waveform} from the process.

\item \textbf{Fix time $(t = t_1)$:}

$X(\xi, t_1) = X(t_1)$ becomes a \textbf{random variable (RV)}.

\item \textbf{Multiple time instances}:

$X(t_1), X(t_2)$ are \textbf{dependent random variables}.
\end{enumerate}

\subsection*{Reasoning}

\begin{itemize}
\item Fixing time $\rightarrow$ randomness remains $\rightarrow$ gives RV.
\item Varying time $\rightarrow$ multiple RVs $\rightarrow$ their \textbf{dependency defines process behavior}.
\end{itemize}

\section*{2. Characterization of Stochastic Processes}

\subsection*{Probability Density Functions (PDFs)}

\textbf{1st Order}

\[
f_X(x, t_1)
\]
\textbf{2nd Order}

\[
f_{X_1 X_2}(x_1, x_2, t_1, t_2)
\]
\textbf{nth Order}

\[
f_{X_1 \dots X_n}(x_1, \dots, x_n, t_1, \dots, t_n)
\]

\subsection*{Reasoning}

\begin{itemize}
\item Higher-order PDFs capture \textbf{more dependencies}.
\item Full knowledge of all orders $\rightarrow$ complete process description.
\end{itemize}

\section*{3. Statistical Moments}

\subsection*{Mean (First Order Moment)}

\textbf{Definition}

\[
\mu_X(t) = E[X(t)]
\]
\textbf{Using PDF}

\[
\mu_X(t) = \int_{-\infty}^{\infty} x \cdot f_X(x,t)\, dx
\]

\subsection*{Reasoning}

\begin{itemize}
\item Expectation gives \textbf{average value at time $(t)$}.
\item Since PDF may depend on time $\rightarrow$ mean is generally \textbf{time-varying}.
\end{itemize}

\subsection*{Autocorrelation Function}

\textbf{Definition}

\[
R_{XX}(t_1, t_2) = E[X(t_1)X^*(t_2)]
\]
\textbf{Using joint PDF}

\[
R_{XX}(t_1,t_2)=\int\int x_1 x_2 f_{X_1X_2}(x_1,x_2,t_1,t_2)\,dx_1dx_2
\]

\subsection*{Reasoning}

\begin{itemize}
\item Measures \textbf{relationship between values at two times}.
\item Important for \textbf{prediction and estimation}.
\end{itemize}

\subsection*{Properties of Autocorrelation}

\textbf{1. Conjugate Symmetry}

\[
R_{XX}(t_1,t_2) = R_{XX}^*(t_2,t_1)
\]
\textbf{Proof (step-by-step from slides)}
\begin{enumerate}
\item Start:
\[
R_{XX}(t_2,t_1) = E[X(t_2)X^*(t_1)]
\]
\item Take conjugate:
\[
(E[X(t_2)X^*(t_1)])^*
\]
\item Rearrange:
\[
= E[X^*(t_2)X(t_1)]
\]
\item Final:
\[
= E[X(t_1)X^*(t_2)] = R_{XX}(t_1,t_2)
\]
\end{enumerate}
\textbf{2. Positive Definiteness}

\[
\sum_{i=1}^n \sum_{j=1}^n a_i^* a_j R_{XX}(t_i,t_j) > 0
\]

\subsection*{Reasoning}

\begin{itemize}
\item Represents \textbf{quadratic form}.
\item Must be positive for any non-zero vector $\rightarrow$ ensures valid correlation structure.
\end{itemize}

\section*{4. Covariance Function}

\subsection*{Definition}

\[
C_{XX}(t_1,t_2) = E[(X(t_1)-\mu_X(t_1))(X^*(t_2)-\mu_X(t_2))]
\]

\subsection*{Relation}

\[
C_{XX}(t_1,t_2) = R_{XX}(t_1,t_2) - \mu_X(t_1)\mu_X(t_2)
\]

\subsection*{Reasoning}

\begin{itemize}
\item Removes mean $\rightarrow$ gives \textbf{pure dependence}.
\end{itemize}

\section*{5. Positive Definite Matrices}

\subsection*{Conditions}

\begin{enumerate}
\item All eigenvalues $> 0$
\item All leading principal minors $> 0$
\end{enumerate}

\subsection*{Example (from slides)}

\begin{itemize}
\item Determinants:
\begin{itemize}
\item $|A_1| = 2 > 0$
\item $|A_2| = 3 > 0$
\item $|A_3| = 4 > 0$
\end{itemize}
\end{itemize}

\subsection*{Conclusion}

$\rightarrow$ Matrix is \textbf{positive definite}

\section*{6. Stationarity}

\subsection*{Concept}

A process is \textbf{stationary} if:

\begin{itemize}
\item Its statistical properties \textbf{do not change under time shift}
\end{itemize}

\subsection*{Strict Sense Stationary (SSS)}

\subsection*{Definition}

\[
f_X(x_1,\dots,x_n; t_1,\dots,t_n) = f_X(x_1,\dots,x_n; t_1+\tau,\dots,t_n+\tau)
\]

\subsection*{Reasoning}

\begin{itemize}
\item Entire \textbf{joint PDF is invariant}
\item Must hold for \textbf{all orders $(n)$} $\rightarrow$ very strong condition
\end{itemize}

\subsection*{Key Result}

\[
R_{XX}(t_1,t_2) = R_{XX}(t_1 - t_2)
\]

\subsection*{Reasoning}

\begin{itemize}
\item Time shift removes dependence on absolute time
\item Only \textbf{time difference matters}
\end{itemize}

\subsection*{Wide Sense Stationary (WSS)}

\subsection*{Conditions}

\begin{enumerate}
\item Mean is constant:
\[
E[X(t)] = \mu_X = \text{constant}
\]

\item Autocorrelation depends only on lag:
\[
R_{XX}(t_1,t_2) = R_{XX}(t_1 - t_2)
\]
\end{enumerate}

\subsection*{WSS vs SSS}

\begin{itemize}
\item SSS $\Rightarrow$ WSS
\item WSS $\nRightarrow$ SSS
\end{itemize}

\subsection*{Special Case (Gaussian Process)}

\begin{itemize}
\item If Gaussian:
\begin{itemize}
\item Mean + covariance fully define process
\item Hence:
\[
\text{WSS} \Rightarrow \text{SSS}
\]
\end{itemize}
\end{itemize}

\section*{7. Example 1: Proving WSS}

\subsection*{Given}

\[
X(t) = A\cos(\lambda t) + B\sin(\lambda t)
\]

Conditions:

\begin{itemize}
\item $E[A]=E[B]=0$
\item $E[A^2]=E[B^2]$
\item $E[AB]=0$
\end{itemize}

\section*{Step 1: Mean}

\[
E[X(t)] = E[A\cos(\lambda t) + B\sin(\lambda t)]
\]

Using linearity:

\[
= \cos(\lambda t)E[A] + \sin(\lambda t)E[B]
\]

Substitute:

\[
= 0
\]

\checkmark\ Constant $\rightarrow$ Condition 1 satisfied

\section*{Step 2: Autocorrelation}

\[
R_{XX}(t_1, t_2) = E[X(t_1)X(t_2)]
\]

Substitute:

\[
= E[(A\cos\lambda t_1 + B\sin\lambda t_1)(A\cos\lambda t_2 + B\sin\lambda t_2)]
\]

Expand:
\begin{itemize}
    \item Cross terms vanish (since $E[AB] = 0$)
    \item Equal variances used
\end{itemize}

\[
= E[A^2](\cos\lambda t_1 \cos\lambda t_2 + \sin\lambda t_1 \sin\lambda t_2)
\]

Use identity:

\[
\cos a \cos b + \sin a \sin b = \cos(a - b)
\]

\[
= E[A^2]\cos\lambda(t_1 - t_2)
\]

\checkmark\ Depends only on $t_1 - t_2$

$\rightarrow$ Process is \textbf{WSS}

\section*{8. Example 2: Random Phase Cosine}

\textbf{Given}

\[
V(t) = \cos(\omega t + \theta)
\]

\begin{itemize}
    \item $\theta \sim \text{Uniform}[-\pi, \pi]$
\end{itemize}

\section*{Step 1: Mean}

\[
E[V(t)] = \frac{1}{2\pi} \int_{-\pi}^{\pi} \cos(\omega t + \theta)\, d\theta
\]

Integrate:

\[
= \frac{1}{2\pi} \left[\sin(\omega t + \theta)\right]_{-\pi}^{\pi}
\]

\[
= 0
\]

\checkmark\ Independent of time

\section*{Step 2: Second Moment}

\[
E[V^2(t)] = \frac{1}{2\pi} \int \cos^2(\omega t + \theta)\, d\theta
\]

Use identity:

\[
\cos^2 x = \frac{1 + \cos(2x)}{2}
\]

Integral gives:

\[
= \frac{1}{2}
\]

\checkmark\ Independent of time

\section*{9. Ergodic Process}

\subsection*{Definition}

A process is \textbf{ergodic} if:

\begin{itemize}
\item \textbf{Time average = Ensemble average}
\end{itemize}

\subsection*{Reasoning}
\begin{itemize}
\item Ideally:
\begin{itemize}
\item Average over many realizations
\end{itemize}
\item Practically:
\begin{itemize}
\item Only \textbf{one realization available}
\end{itemize}
\end{itemize}

If:

\begin{itemize}
\item Single realization gives correct average\\
     $\rightarrow$ Process is ergodic
\end{itemize}

\subsection*{Time Averages}

\textbf{Mean}

\[
\langle X(t)\rangle = \lim_{T\to\infty}\frac{1}{T}\int_{-T}^{T}x(t)\,dt
\]
\textbf{Autocorrelation}

\[
R_{XX}(\tau) = \langle x(t)x(t+\tau)\rangle
\]

\subsection*{Key Result}

\begin{itemize}
\item If WSS + ergodic:
\begin{itemize}
\item Time averages = statistical averages
\end{itemize}
\end{itemize}

\section*{Final Summary}

\begin{itemize}
\item \textbf{Stochastic process} $\rightarrow$ collection of random waveforms
\item \textbf{Mean \& autocorrelation} $\rightarrow$ key statistical tools
\item \textbf{SSS} $\rightarrow$ full distribution invariant
\item \textbf{WSS} $\rightarrow$ only mean + autocorrelation invariant
\item \textbf{Ergodic} $\rightarrow$ time averages = ensemble averages
\end{itemize}

\end{document}